\section{Introduction}
In daily life, we frequently encounter problems where input data is not a deterministic single value, but is only known to belong to some set of possible scenarios. We consider the following two introductory examples:

\begin{example}
	A coffee shop $A$ serves a daily customer count belonging to set $K = \set{10, 15, 30}$ (customers). The average spending per customer belongs to set $C = \set{10, 20, 50}$ (thousand VND). The set of all possible daily revenue levels for coffee shop $A$ is determined by:
	\[
	D = \set{kc : k \in K, \, c \in C} = \set{100, 150, 200, 300, 500, 600, 750, 1500} \text{ (thousand VND)}
	\]
\end{example}

\begin{example}
	Mr. $A$ plans to travel from Hanoi to Ho Chi Minh City via 2 legs:
	\begin{itemize}
		\item Leg $A_1$: Possible distance $S_1 \in \set{200, 240, 300}$ (km), travel speed $V_1 \in \set{40, 50}$ (km/h).
		\item Leg $A_2$: Possible distance $S_2 \in \set{400, 500, 600}$ (km), travel speed $V_2 \in \set{80, 100, 120}$ (km/h).
	\end{itemize}
	The possible travel times on each leg are given respectively by the sets:
	\[
	T_1 = \set{\frac{s_1}{v_1} : s_1 \in S_1, \, v_1 \in V_1} = \set{4; \, 4.8; \, 5; \, 6; \, 7.5} \text{ (hours)}
	\]
	\[
	T_2 = \set{\frac{s_2}{v_2} : s_2 \in S_2, \, v_2 \in V_2} = \set{3.33; \, 4; \, 4.17; \, 5; \, 6; \, 6.25; \, 7.5} \text{ (hours)}
	\]
	Thus, the total possible travel time for the whole journey belongs to the set:
	\[
	T = \set{t_1 + t_2 : t_1 \in T_1, \, t_2 \in T_2} \text{ (hours)}
	\]
	From this, we can derive characteristic laws of the result set $T$, such as minimum value $\min(T)$, maximum value $\max(T)$, mean of $T$, or determine whether $T$ satisfies a certain condition function or criterion $f(T)$. Grasping this complete structural range forms the core foundation for solving decision-making problems under data uncertainty.
\end{example}

To automate and systemize direct computations over these sets of possible scenarios, we construct \textbf{Uncertain Arithmetic} — an extension tool of classical arithmetic for quantities exhibiting uncertainty.
\section{Uncertain Number}
\begin{definition}[Uncertain Number over Field $\mathbb{K}$]
	Let $\mathbb{K}$ be a field; an uncertain number $A$ is a quantity whose value belongs to a set $A \subseteq \mathbb{K}$. The set of all uncertain numbers over field $\mathbb{K}$ is denoted by $\tsbd{K}$. If $A =\s{a}$, we identify $A$ with element $a$. If $A =\emptyset$, we say $A$ is undefined. The number $\emptyset \in \tsbd{K}$ is called the null number.
\end{definition}
\begin{definition}[Transfer Operator]
	To resolve notation ambiguity between sets and uncertain numbers in certain cases, we define the following two transfer operators:
	\begin{enumerate}
		\item Given $A \in \tsbd{K}$, we denote $A_s$ as the set containing all possible values of $A$.
		\item Given $A$ as a set, the operator transforming $A$ into an uncertain number is denoted by $A_u \in \tsbd{K}$.
	\end{enumerate}
\end{definition}
\begin{remark}
	In later sections of this document, we will prove $f(A_s) \neq f(A)$ when $A$ is an uncertain number. This is why uncertain numbers are denoted by $\{\circ\}_u$ rather than standard set notation.
\end{remark}
\begin{remark}
	Throughout this document, we adopt the conventions:
	\begin{enumerate}[leftmargin=*]
		\item Single-valued numbers are denoted by lowercase letters such as $x, y, \dots$.
		\item Uncertain numbers are denoted by uppercase letters such as $X, Y, \dots$.
		\item Sets of uncertain numbers are denoted by calligraphic uppercase letters such as $\mathcal{X}, \mathcal{Y}, \dots$.
		\item Arithmetic spaces are notation conventions to avoid constructing multiple distinct symbols for the same operation within each arithmetic space. The following is the list of arithmetic space notations:
		\begin{enumerate}
			\item $\kgdd{\circ}$: single-point arithmetic space.
			\item $\kgddmr{\circ}$: extended single-point arithmetic space.
			\item $\kgdt{\circ}$: Minkowski arithmetic space.
			\item $\kgdtmr{\circ}$: extended Minkowski arithmetic space.
		\end{enumerate}
	\end{enumerate}
\end{remark}
\begin{definition}[Dimension]
	The finite dimension of a finite-element uncertain number $A \in \tsbd{K}$ equals the number of elements in $A$, denoted by $\dimU(A)$. If $A$ contains infinitely many elements, we say the dimension of $A$ is infinite, written as $\dimU(A) = \infty$.
\end{definition}
\begin{definition}[Canonical Form]
	Let $A \in \tsbd{K}$. We call $f_A : d_A \to \mathbb{K}$ a generating function of uncertain number $A$ if there exists an index set $d_A \subseteq \mathbb{K}$ satisfying: \[d_A = \begin{cases*}
		\set{1..n} \subseteq \mathbb{N^*} \text{ if } \dimU(A)=n \\
		\mathbb{N^*} \text{ if } A_s \text{ is a countably infinite set  }\\
		[0,1] \text{ if } A_s \text{ is an uncountably infinite set.}
	\end{cases*}\] such that $\bd{f_A(d_A)}= A$.
	In this case, we call $(f_A,d_A)$ a canonical form of $A$ and denote $A \sim (f_A,d_A)$.
\end{definition}
\begin{theorem}[Equivalent Canonical Form]
	Let $A \in \tsbd{K}$, and $(f_A,d_A)$ be a canonical form of $A$. Then there exists $(f'_A,d_A)$ which is also a canonical form of $A$.
\end{theorem}
\par{Next, related to solving uncertain equations, we arrive at the following definition:}
\begin{definition}[Uncertain Number of Order $n$]
	Let $A \in \tsbd{K}$, $n$ be a positive integer, we denote $\tsbdcapn{K}{n}$ as the set of order-$n$ uncertain numbers over field $\mathbb{K}$, given by: 
	\begin{enumerate}
		\item For $n = 2$, the set of order-2 uncertain numbers is defined by:
		\[
		\tsbdcapn{K}{2} := \s{A : A \in \tsbd{K}}
		\]
		\item For all $n > 2$, the set of order-$n$ uncertain numbers is defined inductively:
		\[
		\tsbdcapn{K}{n} := \s{A : (\exists p \in \set{1, \dots, n-1})(A \in \tsbdcapn{K}{p})}
		\]
		where by convention $\tsbdcapn{K}{1} \equiv \tsbd{K}$.
		\item For all $n \geq 2$ and $\s{B} \in \tsbdcapn{K}{n}$, hierarchical identification convention is defined by:
		\[
		\s{B} \equiv \begin{cases}
			B & \text{with } B \in \tsbdcapn{K}{n-1} \quad (\text{if } n > 2) \\
			B & \text{with } B \in \tsbd{K} \quad (\text{if } n = 2)
		\end{cases}
		\]
	\end{enumerate}
\end{definition}
\begin{definition}[Uncertain Radius]
	Let $A$ be an uncertain number, the uncertain radius of $A$ is denoted by $\rad(A)$ and $$\rad(A) = \frac{1}{2} \cdot \sup \set{x_1 - x_2 : x_1, x_2 \in A}$$
\end{definition}
\begin{corollary}
	An uncertain number $A = \s{a}$ with $a \in \mathbb{K}$ has an uncertain radius $\rad(A) = 0$.
\end{corollary}
\begin{definition}[Distance]
	Let $\mathbb{K}$ be a normed field, $A, B \in \tsbd{K}$. The distance between two uncertain numbers $A$ and $B$ is denoted by $d_H(A,B)$, given by
	$$
	d_H(A,B) := \max\left\{\sup_{a \in A}{d(a,B),\sup_{b \in B}{d(b,A)}}\right\}
	$$ Where
	$$d(a,A) := \inf_{b\in A}{|a-b|}$$
\end{definition}
\begin{definition}[Weak Binary Relation]
	Consider $\mathcal{R}$ as a binary relation on field $\mathbb{K}$ and two uncertain numbers $A \sim (f_A, d_A)$, $B \sim (f_B, d_B)$ in $\mathbf{U}(\mathbb{K})$. 
	
	The weak binary relation $A \mathcal{R} B$ is a proposition evaluated on scenario space $\Omega = d_A \times d_B$ with atomic truth function:
	\[
	v(A \mathcal{R} B, \omega) = \begin{cases} 
		1 & \text{if } f_A(t_1) \mathcal{R} f_B(t_2) \text{ is true}, \\
		0 & \text{if } f_A(t_1) \mathcal{R} f_B(t_2) \text{ is false},
	\end{cases} \quad \text{with } \omega = (t_1, t_2) \in \Omega
	\]
	Then, the truth value of relation $A \mathcal{R} B$ is determined by:
	\[
	\ddclmd(A \mathcal{R} B) := \frac{\lambda\big(\{(t_1, t_2) \in d_A \times d_B : f_A(t_1) \mathcal{R} f_B(t_2)\}\big)}{\lambda(d_A \times d_B)}
	\]
\end{definition}
\begin{example}
	Given $A = \s{1,2}$ as an uncertain number, we have $$\gtcl{A\leq A}{0.75}$$ $$\gtcl{A>A}{0.25}$$
	We remark: $\ddclmd(A\leq A) + \ddclmd(A>A) = 1$. Although $A=A$, their sizes differ probabilistically/scenariowise.
\end{example}
\begin{example}
	Given $A = \s{1,2}$ and $B = \s{3} = 3$, then $A \leq B$.
\end{example}
\begin{theorem}[Balance]
	Applying the definition of weak binary relation with $\mathcal{R}$ as $\leq$, we have the following theorem:
	for $X \in \tsbd{R}$ where $\dimU(X) = \infty$, then: 
	$$\ddclmd(X\leq X) = \frac{1}{2}$$
\end{theorem}
\begin{proof}
	Suppose $\taphop(X)$ has $n$ distinct elements. Total possible pairs $(a,b)$ is $n^2$. The number of pairs where $a<b$ is $n(n-1)/2$. The number of pairs where $a=b$ is $n$. Therefore, the number of pairs $a\leq b$ is $$\frac{n(n-1)}{2} + n = \frac{n^2+n}{2}$$ 
	Thus, the truth value is $$\ddclmd(X\leq X ) = \frac{n^2+n}{2n^2} = \frac{1}{2} + \frac{1}{2n}$$
	Letting $n \rightarrow \infty$ completes the proof.
\end{proof}
\begin{definition}[Set Arithmetic Operations]
	Let $*$ be an arbitrary set operation, $A, B \in \mathbf{U}(\mathbb{K})$, we define:
	\[
	A * B := (A_s * B_s)_u
	\]
	Where $A_s * B_s$ on the right hand side represents the set operation.
\end{definition}
\begin{example}
	Consider the problem: Person $A$ has an uncertain number of apples $A=\s{1,2}$. $A$ intends to share these apples with room $B$, where room $B$ has an uncertain number of people $B = \s{1,0}$. After $A$ divides the apples, each person in room $B$ receives a quantity of apples equal to \[\frac{A}{B} = \s{1,2}\] Why does $\frac{1}{0}$ not appear in the result set? Because if $A$ sees no people in the room, $A$ does not perform division, nor can $A$ divide apples to each person in room $B$ with $+\infty$ apples. This division does not happen rather than yielding no valid outcome. Hence we define binary operations on the set of uncertain numbers as follows:
\end{example}
\begin{definition}[Single-Point Arithmetic Space]
	Let $A \in \tsbd{K}$ be an uncertain set. We define a \textbf{single-point uncertain function} as a mapping $f$ acting on variable $A$, such that during each evaluation pass, variable $A$ receives a unique representative point value $x \in A$ at all occurrences of $A$ in expression $f$.
	
	The image set of the corresponding point-wise uncertain function, denoted $\kgdd{f}$, is determined by:
	\[
	\kgdd{f}(A) := \s{f(x) \;:\; x \in A}
	\]
	Operations within $(\circ)_1$ are called operations on the single-point arithmetic space.
\end{definition}
\begin{example}
	Let $f : \mathcal{X} \to \mathcal{Y}$ be given by:
	\[f(X) = (X^2+5X+6)_1\] as a single-point uncertain function.
\end{example}
\begin{remark}
	This definition ensures algebraic consistency via two core mechanisms:
	\begin{itemize}
		\item \textbf{Internal Point Identity:} Every symbol of variable $A$ appearing in expression $f$ (regardless of position or operation) is bound to receive the same value $x \in A$ during the same mapping pass.
		\item \textbf{Elimination of Quantity Expansion:} Thanks to unified point assignment, classical algebraic properties are fully preserved, such as $\kgdd{x - x}(A) = \s{0}$ and $\kgdd{x + x}(A) = 2A$.
	\end{itemize}
\end{remark}

\begin{example}
	Consider single-variable function $f(A) = A^2 - 3A + 2$ with $A \in \tsbd{\mathbb{R}}$. The image set of this single-point uncertain function is defined by:
	\[
	\kgdd{f}(A) = \s{x^2 - 3x + 2 \;:\; x \in A}
	\]
	During each evaluation $x \in A$, both positions $A^2$ and $-3A$ uniformly receive the same value $x$.
\end{example}
\begin{definition}[Minkowski Arithmetic Space]
	Let $\mathbb{K}$ be a field, with * as any binary operation on $\mathbb{K}$, we define arithmetic binary operation * on $\tsbd{K}$:
	For $A,B \in \tsbd{K}$, then $$(A * B)_m = \s{a*b : a \in A , b\in B , a*b \in \mathbb{K}}$$
	\noindent If $A * B =\emptyset$, we say $A * B$ is undefined. In writing, we abbreviate $(A*B)_m$ as $A*B$ and refer to operations in $(\circ)_m$ as operations executed on Minkowski arithmetic space.
\end{definition}
\begin{definition}[Iterated Operations]
	Let $*$ be any binary operation on the set of uncertain numbers, $A$ be an uncertain number, $n$ be a natural number greater than 1. We define: $$A^{n}_{*} := \begin{cases}
		A*A & n=2 \\
		\left(A^{n-1}_*\right) * A & n > 2
	\end{cases}
	$$\
\end{definition}
\begin{example}
	\begin{enumerate}
		\item If * is addition on uncertain numbers, we have $$A^n_+ = A+A+... \text{($n$ times)}$$
		\item If * is multiplication on uncertain numbers, we have $$A^n_* = A.A .... \text{($n$ times)}$$
		\item If * is min operation on uncertain numbers, we have
		$$A^n_{\min} = \begin{cases}
			\min(A,A) & n=2 \\
			\min\left(A^{n-1}_{\min}, A\right) & n > 2
		\end{cases} $$
		It is easy to see $A^n_{\min} = A, \forall n \geq 2$
	\end{enumerate}
\end{example}
\begin{example}
	Let $2_*$ be an iterated operation where $*$ is an arbitrary binary operation on the set of real numbers, $A$ is an uncertain number, we have:
	\[A^2_* = \bigcup_{a \in A} A * a \]
\end{example}
\par{Operations written inside notation $\kgdtmr{\circ}$ are called operations on extended Minkowski arithmetic space. One of them is defined as follows:}
\begin{definition}
	Let $p,q \in \mathbb{N}^*$.
	We define \[\kgdtmr{\frac{p}{q}A} := \s{X \in \tsbd{K}: X^q_+ = A^p_+} \]
	If $q = 1$, we set: \[\kgdtmr{pA}:= \kgdtmr{\frac{p}{1} A} \]
\end{definition}
\begin{example}
	$\kgdtmr{\frac{1}{2} \s{2,3,4}} = \s{\s{1,2}} = \s{1,2}$
\end{example}
\begin{definition}
	Let $p,q \in \mathbb{N^*}$.
	We define \[\kgdtmr{A^{\frac{p}{q}}} := \s{X \in \tsbd{K}: X^q_* = A^p_*} \]
	If $p = 1$, we set: \[\kgdtmr{\sqrt[q]{A}} := \kgdtmr{A^{\frac{1}{q}}}\]
\end{definition}
\begin{example}
	To help readers avoid confusing arithmetic spaces in this study, we state the following problems:
	\begin{enumerate}
		\item Mr. $A$ plans to build a square courtyard with side length belonging to set $\set{1,2}$. Since after choosing 1 side, the remaining side length must equal that chosen number, the area of the courtyard is: \[S_1 = \kgdd{\s{1,2}^2} = \s{1,4}\]
		\item Mr. $B$ plans to build a courtyard but assigns two decision-makers for side lengths; each chooses a length from set $\set{1,2}$. Then the area of Mr. $B$'s courtyard is:
		\[S_2 = \kgdtmr{\s{1,2}^2} =\s{1,2,4}\]
	\end{enumerate}
	Thus, total courtyard area for Mr. $A$ and Mr. $B$ combined is:\[S= S_1 + S_2 = \s{2, 3, 5, 6, 8}\]
	Let $X \in \tsbd{K}$ be an uncertain number denoting possible choices for building courtyards of Mr. $A$ and $B$, the total area function equals:
	\[f(X) = \kgdd{X^2} + \kgdtmr{X^2}\] which is an uncertain function.
\end{example}
\begin{theorem}[Canonical Form on Minkowski Arithmetic Space]
	Let $A,B \in \tsbd{K}$, $*$ be any binary operation in Minkowski arithmetic space, and $(f_A,d_A), (f_B,d_B)$ be canonical forms of $A$ and $B$ respectively. Then the canonical form of $\kgdt{A * B}$ is determined by generating function:
	\[f_{A*B} : d_A \times d_B \to \mathbb{K}\]
	With \[f_{A*B}(x_1,x_2) = f_A(x_1) * f_B(x_2)\]
	Therefore: \[\kgdt{A*B} = \s{f_A(x_1)*f_B(x_2) : (x_1,x_2) \in d_A \times d_B}\]
\end{theorem}
\par{\noindent Next we present an example computing with large numbers:}
\begin{example}
	Consider two uncertain numbers $A = \{1, 3\}_u$ and $B = \{10, 20\}_u$ over field $\mathbb{R}$.
	
	Choose index sets $d_A = d_B = \{1, 2\}$. By Lagrange interpolation, we construct polynomial generating functions $f_A: d_A \to \mathbb{R}$ and $f_B: d_B \to \mathbb{R}$ as follows:
	$$f_A(x_1) = 1 \cdot \frac{x_1 - 2}{1 - 2} + 3 \cdot \frac{x_1 - 1}{2 - 1} = 2x_1 - 1$$
	$$f_B(x_2) = 10 \cdot \frac{x_2 - 2}{1 - 2} + 20 \cdot \frac{x_2 - 1}{2 - 1} = 10x_2$$
	
	Then $(f_A, d_A)$ and $(f_B, d_B)$ are polynomial canonical forms of $A$ and $B$ respectively.
	
	By Theorem 2.3, the generating function of $(A + B)_m$ on domain $d_A \times d_B$ is defined by two-variable polynomial:
	$$f_{A+B}(x_1, x_2) = f_A(x_1) + f_B(x_2) = 2x_1 + 10x_2 - 1$$
	
	Evaluating polynomial $f_{A+B}$ at each point in $d_A \times d_B$:
	\begin{align*}
		f_{A+B}(1, 1) &= 2(1) + 10(1) - 1 = 11, \quad f_{A+B}(1, 2) = 2(1) + 10(2) - 1 = 21 \\
		f_{A+B}(2, 1) &= 2(2) + 10(1) - 1 = 13, \quad f_{A+B}(2, 2) = 2(2) + 10(2) - 1 = 23
	\end{align*}
	
	Therefore:
	$$ A + B = \{11, 13, 21, 23\}_u$$
	From which if we set $P = A+B$, we compute \[\kgdtmr{P^{100}} = \s{\prod_{i=1}^{100}{(2x_i + 10 y_i - 1)}: (x_i,y_i) \in d_A \times d_B, i \in \{1..100\}}
	\]
\end{example}
\begin{proposition}
	Let $A \sim (f_A,d_A), B \sim (f_B,d_B)$ in $\tsbd{K}$, consider $a \in A$, then $a \in B$ if and only if \[
	\exists (x,y) \in d_A \times d_B : a=f_A(x)=f_B(y)
	\]
\end{proposition}
\begin{remark}
	This proposition transforms the essence of element checking $a \in B$ (or subset inclusion $A \subseteq B$) from cumbersome data lookups into finding roots of algebraic equation $f_B(y) = f_A(x)$ over index domain $d_A \times d_B$. This completely eliminates dependence on data size when sets expand.
\end{remark}
\begin{proposition}[Canonical Form of Intersection]
	\label{prop:dang_chuan_tac_phep_giao}
	Let $A\sim (f_A,d_A), B \sim (f_B,d_B)$ in $\tsbd{K}$. Then \[
		C = A \cap B
	\]
	is an uncertain number with canonical form $(f_{A\cap B},d_{A\cap B})$ given by:
	\begin{align*}
		d_{A\cap B} &= \set{(x,y) \in d_A \times d_B : f_A(x) = f_B(y)} \\
		f_{A\cap B}(z) &= f_A(x) = f_B(y), z=(x,y) \in d_{A\cap B}
	\end{align*}
\end{proposition}
\begin{example}
	Consider two uncertain numbers $A, B \in \mathbf{U}(\mathbb{R})$ with canonical forms given by:
	\begin{itemize}
		\item $A \sim (f_A, d_A)$ with $f_A(x) = 2x + 1$ on $d_A = \{1, 2, 3, 4\} \Rightarrow A = \{3, 5, 7, 9\}_u$.
		\item $B \sim (f_B, d_B)$ with $f_B(y) = 3y - 1$ on $d_B = \{1, 2, 3, 4\} \Rightarrow B = \{2, 5, 8, 11\}_u$.
	\end{itemize}
	
	According to Proposition \ref{prop:dang_chuan_tac_phep_giao}, we find index domain $d_{A \cap B}$ by solving algebraic equation on $d_A \times d_B$:
	\[
	f_A(x) = f_B(y) \iff 2x + 1 = 3y - 1 \iff 2x - 3y = -2
	\]
	Testing pairs $(x, y) \in \{1, 2, 3, 4\} \times \{1, 2, 3, 4\}$, the equation has unique solution $x = 2, y = 2$. Therefore:
	\[
	d_{A \cap B} = \{(2, 2)\}
	\]
	The corresponding generating function yields value:
	\[
	f_{A \cap B}(2, 2) = f_A(2) = 2(2) + 1 = 5
	\]
	Thus canonical form $(f_{A \cap B}, d_{A \cap B})$ accurately identifies $A \cap B = \{5\}_u$.
\end{example}
\begin{proposition}[Canonical Form of Union]
	Let $A \sim (f_A, d_A)$ and $B \sim (f_B, d_B)$ in $\mathbf{U}(\mathbb{K})$. Canonical form $(f_{A \cup B}, d_{A \cup B})$ of uncertain number $A \cup B$ is defined by:
	\[
	\begin{cases}
		d_{A \cup B} = (d_A \times \{0\}) \cup (d_B \times \{1\}) \\[1ex]
		f_{A \cup B}(z) = \begin{cases} 
			f_A(x), & \text{if } z = (x, 0) \in d_A \times \{0\} \\[1ex]
			f_B(y), & \text{if } z = (y, 1) \in d_B \times \{1\}
		\end{cases}
	\end{cases}
	\]
\end{proposition}
\begin{theorem}[Undefinedness]
	Let $*$ be any binary operation on $\tsbd{K}$, $A \in \tsbd{K}$. Then we have \[A * \emptyset = \emptyset * A= \emptyset\]
	That is, any uncertain arithmetic binary operation involving the null number yields the null number (undefined).
\end{theorem}
\begin{example}
	With operation * as +, we have the definition of addition on $\tsbd{K}$:
	$$A + B = \s{a+b : a \in A , b \in B}, \text{ with } A,B \in \tsbd{K}.$$
\end{example}
\begin{example}
	$$\frac{\s{1,2,3}}{\s{2,3}} = \s{1,\frac{1}{3},\frac{1}{2},\frac{2}{3},\frac{3}{2}}$$
	$$\frac{\s{1,2,3}}{\s{1,0}} = \s{1,2,3}$$
	$$\frac{\s{1,2,0}}{0} \text{ is undefined}$$
	$$\frac{0}{0} \text{ is undefined}$$
\end{example}
\begin{md}[Distributive Property of Operation $*$ with Union Operation]
	Let $*$ be a binary operation on uncertain numbers, $A$ and $B$ be uncertain numbers, we have
	\[A*B = \bigcup_{b \in B} {A * b} \]
\end{md}
\begin{theorem}
	Let $A \in \tsbd{K}$, we set \begin{enumerate}
		\item $Id_+=\set{X-X: X \in \tsbd{K}}$
		\item $Id_*=\set{X/X: X \in \tsbd{K}}$
	\end{enumerate}
	From which we have: 
	\begin{align*}
		(\forall e_+ \in Id_+)(A + e_+ &= A - e_+) \\
		(\forall e_* \in Id_*)(A \cdot e_* &= A / e_*)
	\end{align*}
\end{theorem}
\begin{theorem}[Inverse]
	Consider $A \in \tsbd{K}, \dimU(A) > 1$, then:
	\begin{enumerate}
		\item There exists no uncertain number B such that $A+B = 0$
		\item There exists no uncertain number B such that $AB = 1$
	\end{enumerate}
\end{theorem}
\begin{theorem}
	Let $A$ and $B$ be two uncertain numbers, we have $$\rad(A+B) \leq \rad(A) +\rad(B)$$
\end{theorem}
\begin{theorem}
	Let $A$ be an uncertain number, $b$ be a single-valued constant. We have $$\rad(A+b) = \rad(A)$$
	This theorem allows an object with uncertain position to preserve its uncertain radius under rigid translation.
\end{theorem}
\begin{theorem}[Realistic Similarity]
	Let $A$ be an uncertain number and $k$ be a single-valued constant, we have: $\rad(k A) = |k| \rad(A)$
\end{theorem}

\begin{theorem}[Identity Theorem with Classical Arithmetic]
	Let $A = \s{a} , B = \s{b}$ and * be a binary operation on $\mathbb{R}$ where $a*b$ is defined, then $$A * B = \s{a*b} = a*b$$
\end{theorem}
\begin{theorem}[Commutativity]
	Let $A,B$ be two uncertain numbers and * be a commutative binary operation on $\mathbb{R}$, we have
	$$A*B = B*A$$
\end{theorem}
\begin{theorem}[Associativity]
	Let $A,B,C$ be uncertain numbers and * be an associative binary operation on $\mathbb{R}$, then:
	\[A*(B*C) = (A*B)*C\]
\end{theorem}
\begin{theorem}[Distributivity]
	Let $K,A,B$ be uncertain numbers, we have:
	\[K(A+B)\ = \bigcup_{k \in K} [k(A+B)] = \bigcup_{k \in K} (kA+kB)\]
	\[(AB)^K\ = \bigcup_{k \in K} \left[(AB)^k \right] = \bigcup_{k \in K} \left(A^kB^k\right) \]
\end{theorem}
\begin{theorem}[Inter-Arithmetic Space Theorem]
Let $a,c \in \mathbb{K}, X \in \tsbd{K}$. We have:
\begin{enumerate}
	\item $\kgdd{aX} = \kgdt{aX}$
\end{enumerate}
\end{theorem}
\section{Weighted Uncertain Numbers}
\begin{definition}[Weight]
	Let $A$ be a set, a function $\ts{A}$ from $A$ to $[0,1]$ satisfying the following properties is called a weight of A: 
	\begin{enumerate}
		\item $$\sum_{a \in A}{\ts{A}(a) } = 1$$
		\item For all $a \in A$, $\ts{A}(a) > 0 $
	\end{enumerate}
\end{definition}
\begin{definition}[Weighted Uncertain Number over Field $\mathbb{K}$]
	A pair $(A,\ts{A})$ with $A \in \tsbd{K}$ and $\ts{A}$ being a weight of $A$ is called a weighted uncertain number over field $\mathbb{K}$. The set of all weighted uncertain numbers over field $\mathbb{K}$ is denoted by $\tsbdts{K}$.
\end{definition}
\begin{example}
	Let $A \in \tsbd{R}$ and $A = \s{a} = a $ with arbitrary weight $\ts{A}$, we always have $A = a$ with $\ts{A}(a) =1$. Such a number $A$ is a real number.
\end{example}
\begin{example}
	Set $A=\s{2k : k \in \mathbb{N}}$ together with weight $\ts{A}(2k) = \frac{1}{2^{k+1}}$ is an uncertain number over real field with weight $\ts{A}$.
\end{example}
\begin{theorem}[Weight of an Operation]
	Let * be a binary operation on $\tsbd{K}$, we define function $$\ts{A*B}(c) := \sum_{a*b=c}{\ts{A}(a)\ts{B}(b)}$$ with $c = a * b \in A * B $, where $\ts{A}$ is a weight of $A$ and $\ts{B}$ is a weight of $B$. Then $\ts{A*B}$ is also a weight of $A *B $.
\end{theorem}
\begin{definition}[Weighted Weak Binary Relation over Field $\mathbb{K}$]
	Let $\mathcal{R}$ be a binary relation on field $\mathbb{K}$. $(A,\ts{A}),(B,\ts{B}) \in \tsbdts{K}$. We define the weighted weak binary relation over field $\mathbb{K}$ as follows:
	\[\ddclmd(A\mathcal{R}B) := \sum_{(a,b) \in A \times B}{\left[\ts{A}(a)\ts{B}(b)I(a,b)\right]}\]
	Where I(a,b) is given by \[
	I(a,b) := \begin{cases}
		1 & \text{if } (a,b) \in \mathcal{R} \\
		0 & \text{if } (a,b) \notin \mathcal{R}
	\end{cases}
	\]
\end{definition}